\section{Introduction}\label{sec:intro}

Biological metaphors have a long history in computing.  Artificial
immune systems detect network intrusions~\cite{hofmeyr2000architecture}.
Swarm intelligence algorithms solve combinatorial
optimization~\cite{bonabeau1999swarm}.  Neural architectures borrow
structural motifs from neuroscience.  In each case, the biological
analogy provides organizing structure for the algorithm---but the
question of whether that structure \emph{earns its complexity} over
simpler alternatives is rarely asked.

The default assumption is that biological inspiration provides
qualitative benefits: more robust error handling, more adaptive
resource management, more distributed coordination.  What is missing
is quantitative comparison.  Does a multi-currency metabolic state
machine actually serve more critical operations under pressure than a
flat budget counter?  Does signal accumulation with temporal decay
actually reduce false alarms compared to majority voting?  Does
Bayesian two-signal stagnation detection actually distinguish
convergence from pathological loops better than cosine similarity?

This paper tests these questions empirically using Operon, a
biologically-inspired AI agent reliability framework that implements
22 biological motifs spanning genome-level configuration through
tissue-level coordination~\cite{banu2026operon}.  We select three
motifs that represent distinct reliability claims, ground each
benchmark in real biological pathway data from KEGG and Reactome,
and compare each biological implementation against both a naive
non-biological alternative and an ablated control.

\paragraph{The structural guarantee hypothesis.}
Our central hypothesis is that biological design earns complexity when
it provides \emph{mechanism-level structural guarantees}---hard
properties enforced by design rather than achieved by optimization.
A metabolic state machine that rejects low-priority operations in
STARVING state provides a structural guarantee: critical operations
\emph{will} be served regardless of load pattern.  A signal
accumulation model with exponential decay provides a structural
guarantee: stale evidence \emph{will} age out regardless of query
timing.  These are mechanism design choices, not algorithmic
improvements.

We contrast this with biological designs that attempt
\emph{information processing}---using the biological metaphor to make
better decisions rather than to enforce structural invariants.  Our
hypothesis is that these do not reliably outperform simpler
alternatives.

\paragraph{Pathway-grounded methodology.}
Each benchmark derives its test scenarios from real biological pathway
data.  The metabolism benchmark uses AMPK signaling pathway parameters
(KEGG hsa04152) for metabolic state thresholds and rate-of-change
sensitivity.  The quorum sensing benchmark uses autoinducer kinetics
from the \emph{V.~fischeri} LuxI/LuxR system (KEGG map02024) for
signal decay half-life and threshold scaling.  The epiplexity
benchmark uses the free energy principle~\cite{friston2010free} to
generate trophic-withdrawal scenarios.  This grounding prevents the
benchmarks from testing against our own assumptions about what failure
patterns look like.

\paragraph{Three-variant comparison.}
Each benchmark runs three experimental conditions:
\begin{enumerate}[nosep]
  \item \textbf{Biological}: the full Operon feature, biologically grounded.
  \item \textbf{Ablated}: the feature disabled entirely (e.g., no monitor, no scaling).
  \item \textbf{Naive}: a simple, reasonable, non-biological alternative
        that a competent engineer would reach for (e.g., cosine similarity,
        majority vote, flat counter).
\end{enumerate}
Ablation proves necessity; the naive alternative proves the biological
\emph{design} matters, not just having any feature at all.

\paragraph{Contributions.}
\begin{enumerate}[nosep]
  \item A benchmark suite testing three biologically-grounded agent
        reliability features against naive alternatives, with 10M+
        data points across 10 seeds.
  \item Two clear empirical wins for biological structural guarantees
        (metabolism, quorum sensing) and one honest negative
        (epiplexity stagnation detection).
  \item Evidence that the distinction between \emph{structural
        guarantees} and \emph{algorithmic sophistication} predicts
        where biological design earns complexity.
\end{enumerate}
